\documentclass[aspectratio=169,10pt]{beamer}
\usetheme{metropolis}
\usepackage{booktabs,amsmath,amssymb}
\title{Q18: LLMs vs.\ Brains}
\subtitle{What the analogy gets right and wrong}
\author{DLS 2026 --- Kukla}
\date{October 2026}
\begin{document}
\maketitle

\begin{frame}{The question}
\begin{block}{Student question (\#06, \#11)}
How similar is LLM ``thinking'' to human thinking? How similar are AI systems to biological brains, and
could we build AI from brain structure?
\end{block}
\end{frame}
\begin{frame}{Separate three claims}
\begin{enumerate}
\item \textbf{Behavioral} similarity (same outputs).
\item \textbf{Representational} similarity (same internal codes).
\item \textbf{Mechanistic} similarity (same process).
\end{enumerate}
\alert{Often conflated; very different evidence.}
\end{frame}
\begin{frame}{Where the analogy holds}
\begin{itemize}
\item Distributed representations + superposition (both).
\item Hierarchical features; model activations \emph{predict} sensory cortex.
\item Learning from experience (adapt connection strengths).
\end{itemize}
\end{frame}
\begin{frame}{Where it breaks}
\begin{itemize}
\item \textbf{Learning rule}: backprop is biologically implausible.
\item \textbf{Substrate}: spiking, stochastic, $\sim$20\,W brain vs clocked GPUs.
\item \textbf{Objective}: next-token text vs embodied, evolved, reward-driven.
\item \textbf{Architecture}: attention/KV-cache has no cortical analog.
\end{itemize}
\end{frame}
\begin{frame}{Build AI from brain structure?}
\begin{itemize}
\item Transfers that worked were \textbf{abstractions} (distributed codes, convolution, attention).
\item Literal spiking/connectome models have \emph{not} matched deep learning on tasks.
\item Flow reverses: deep nets are now a leading \emph{model} of sensory cortex.
\end{itemize}
\alert{``Copy the connectome and run it'' is not a proven route to capability.}
\end{frame}
\begin{frame}{Takeaway}
Shared: distributed, hierarchical, experience-shaped representations.\\
Different: learning rule, substrate, objective, architecture.\\[1em]
\alert{Behavioral resemblance $\ne$ mechanistic identity.}
\end{frame}

\end{document}
